\documentclass[11pt]{article}
\usepackage[margin=1in]{geometry}
\usepackage{amsmath,amssymb}
\usepackage{booktabs}
\usepackage{enumitem}
\usepackage{hyperref}
\newcommand{\bs}[1]{\boldsymbol{#1}}
\newcommand{\code}[1]{\texttt{#1}}
\newcommand{\dd}{\mathrm{d}}
\emergencystretch=3em

\title{Anisotropic Energetic Mesh Smoothing with a Metric Tensor Field in Norma}
\author{Alejandro Mota}
\date{September 2026}

\begin{document}
\maketitle

\begin{abstract}
Energetic mesh smoothing (EMS) in Norma minimizes a hyperelastic energy of the
mesh, with an ideal reference element per tetrahedron set by a scalar size.
This note records the extension of that ideal element to a metric tensor,
which prescribes three principal sizes and their directions, and the change
in the energy that makes the orientation of the elements part of the
objective: the energy is evaluated on the deformation gradient measured in
the metric, and the stress and tangent are pulled back to the reference
configuration. The formulation, its properties, the treatment of a metric
that varies in space, the implementation, the regression tests, and two
examples are described.
\end{abstract}

\section{Formulation}
\label{sec:formulation}

\paragraph{The isotropic smoothing energy.} For each tetrahedron the ideal
reference element $\Omega_{\bs X}$ is a regular tetrahedron of edge length
$h$, the current element $\Omega_{\bs x}$ is the element of the mesh, and
$\bs F = \partial \bs x / \partial \bs X$ is the deformation gradient of the
map between them. The smoothing energy of the mesh is
$\Phi = \sum_e V_{\bs X}^e\, W(\bs F_e)$, with $W$ an isotropic
hyperelastic energy density (the Seth--Hill family in Norma) and
$V^e_{\bs X}$ the volume of the ideal element. The nodal positions are the
unknowns and the ideal elements are fixed during a solve. Because $W$ is
isotropic, it depends on $\bs F$ only through the invariants of
$\bs C = \bs F^T \bs F$, and $W(\bs Q \bs C \bs Q^T) = W(\bs C)$ for every
rotation $\bs Q$: rotating the ideal element changes nothing. This is
harmless for a regular ideal element, which has no orientation, and fatal
for an anisotropic one: an elongated element has zero energy in every
orientation, so the energy cannot ask the elongation to point in a
prescribed direction.

\paragraph{The metric tensor.} An anisotropic target is a symmetric
positive definite metric tensor
\begin{equation}
  \bs M \;=\; \bs R\,\operatorname{diag}\!\left(\frac{1}{h_1^2},
  \frac{1}{h_2^2}, \frac{1}{h_3^2}\right) \bs R^T,
  \label{eq:metric}
\end{equation}
with principal sizes $h_i$ and a rotation $\bs R = \exp(\widehat{\bs v})$,
given by a rotation vector $\bs v$, whose columns are the principal
directions. The length of a vector $\bs e$ measured in the metric is
$\sqrt{\bs e \cdot \bs M \bs e}$, so an edge of length $h_i$ along the
$i$-th principal direction has unit length in the metric. Any factor
$\bs F_M$ with $\bs F_M^T \bs F_M = \bs M$ maps physical vectors to
metric-space vectors; Norma uses
\begin{equation}
  \bs F_M \;=\; \operatorname{diag}(1/h_1, 1/h_2, 1/h_3)\, \bs R^T,
  \qquad
  \bs F_M^{-1} \;=\; \bs R\, \operatorname{diag}(h_1, h_2, h_3).
  \label{eq:factor}
\end{equation}
The ideal element is the unit regular tetrahedron $\Omega_{\bs Y}$ mapped by
$\bs F_M^{-1}$, that is, scaled by $h_i$ along the global axes and then
rotated by $\bs R$: $\bs X = \bs F_M^{-1} \bs Y$. Its volume is
$h_1 h_2 h_3\, V_{\mathrm{unit}}$ with $V_{\mathrm{unit}} = 1/(6\sqrt 2)$.

\paragraph{The energy in the metric space.} The current element is also
mapped into the metric space, $\bs y = \bs F_M \bs x$, and the energy is
evaluated on the deformation gradient between the unit ideal element and
the metric image of the current element,
\begin{equation}
  \bs F_U \;=\; \frac{\partial \bs y}{\partial \bs Y}
  \;=\; \bs F_M\, \bs F\, \bs F_M^{-1},
  \qquad
  \Phi \;=\; \sum_e V^e_{\bs X}\, W(\bs F_{U,e}).
  \label{eq:FU}
\end{equation}
An element has zero energy if and only if its metric image is a unit
regular tetrahedron, that is, if it is a unit element for $\bs M$ in the
sense of metric-based meshing. If the current element is the ideal element
rotated by $\bs Q$, then $\bs F = \bs Q$ and
$\bs C_U = \bs F_M^{-T} \bs Q^T \bs M \bs Q \bs F_M^{-1}$, which equals the
identity only when $\bs Q^T \bs M \bs Q = \bs M$: the orientation is
penalized, except for rotations that leave the metric unchanged, as it
should be. With equal sizes $\bs F_M = \bs I / h$, $\bs F_U = \bs F$, and
the isotropic energy is recovered exactly.

\paragraph{Stress and tangent.} With the ideal volume as the weight,
$\Phi_e = V^e_{\bs X} W(\bs F_M \bs F \bs F_M^{-1})$, and the chain rule
gives the first Piola--Kirchhoff stress and the tangent in the reference
configuration from those of the isotropic model evaluated at $\bs F_U$,
\begin{equation}
  \bs P \;=\; \bs F_M^T\, \bs P_U\, \bs F_M^{-T},
  \qquad
  A_{iJkL} \;=\; (F_M)_{pi}\,(F_M^{-1})_{Jq}\, A^U_{pqrs}\,
  (F_M)_{rk}\,(F_M^{-1})_{Ls},
  \label{eq:pullback}
\end{equation}
where $\bs P_U = \partial W/\partial \bs F_U$ and
$\bs A^U = \partial \bs P_U / \partial \bs F_U$. Nothing else in the
assembly changes. The same result follows from the nodal chain rule: with
$\bs y_a = \bs F_M \bs x_a$, the nodal forces are
$\bs f_{\bs x} = \bs F_M^T \bs f_{\bs y}$ and the nodal stiffness blocks
$\bs K_{\bs x} = \bs F_M^T \bs K_{\bs y} \bs F_M$.

\paragraph{Properties.}
\begin{enumerate}[itemsep=2pt]
\item The energy does not depend on the factor chosen for $\bs M$: any
  other factor is $\bs Q \bs F_M$ with $\bs Q$ a rotation, which enters
  $\bs F_U$ as a superimposed rigid rotation and leaves the isotropic
  $W$ unchanged. The eigenvector signs and the in-plane phase of $\bs R$
  for repeated sizes are therefore irrelevant to the energy.
\item The orientation of the unit tetrahedron $\Omega_{\bs Y}$ is
  irrelevant for the same reason: every rotation of the unit tetrahedron
  mapped by $\bs F_M^{-1}$ is a zero-energy element. The smoother does not
  prefer one packing of anisotropic elements over another; the mesh
  topology decides.
\item $\det \bs F_U = \det \bs F$, so the volumetric part of the energy
  is unchanged by the metric transformation: it compares the element
  volume with $h_1 h_2 h_3 V_{\mathrm{unit}}$.
\end{enumerate}

\section{A metric that varies in space}
\label{sec:spatial}

The metric is evaluated once per element at the centroid of the element in
the original mesh, as the isotropic size field is, and held fixed during a
solve. The energy is then an exact function of the nodal positions and the
assembled force is its exact gradient, which the line search and the
convergence tests of the solver require. If the metric followed the
current centroid, $\bs F_M$ would depend on the unknowns and the exact
gradient would contain the derivative of $\bs F_M$, which involves the
derivative of the exponential map of the rotation vector; omitting it
would make the force inconsistent with the energy. The metric is therefore
Lagrangian: an element keeps the target it was assigned where it started.
For smoothing, where nodes move a fraction of an element, this is a
second-order effect. For large motions, or when smoothing alternates with
topological operations, the targets are refreshed at phase boundaries
from the field at the current positions, which is the re-baselining of
the adaptivity note (\code{docs/notes/ems-adaptivity}); each phase then
minimizes a fixed functional.

The metric takes the pseudo-time $t$ as an argument, so a target can be
continued from the sizes of the original mesh to the desired ones over the
pseudo-time steps of the quasi-static solve. Because the sizes and the
rotation vector are given separately, a linear path in $\log h_i$ and in
$\bs v$ interpolates the sizes geometrically and stays a rotation, which is
the interpolation the metric literature recommends.

\paragraph{The volume floor.} Smoothing cannot create elements. The
isotropic \code{size field} rule therefore takes the larger of the size
field and the edge length of the regular tetrahedron with the volume of
the original element, so that no element is asked to shrink below what
the mesh topology can accommodate; without it the size field produced
slivers. The rule \code{metric field} applies the same floor to the
metric: when $h_1 h_2 h_3 V_{\mathrm{unit}}$ is smaller than the volume of
the original element, the three sizes are scaled by the same factor until
the volumes match, so the shape and the orientation of the target are
preserved. The rule \code{metric field unrestricted} omits the floor.

\section{Implementation}
\label{sec:implementation}

\begin{itemize}[itemsep=2pt]
\item \code{MetricField} (\code{src/model\_types.jl}): the compiled
  size and rotation-vector expressions in $t, x, y, z$ and the
  \code{restricted} flag; stored on the model as \code{metric\_field}.
\item \code{create\_metric\_field} (\code{src/model.jl}): reads the
  \code{metric field} block under \code{model} (\code{sizes}, three
  expressions; \code{rotation vector}, three expressions, optional) with
  the same Symbolics pipeline as the size field.
\item \code{create\_metric\_reference}: evaluates the field at the
  reference centroid, checks that the sizes are finite and positive,
  applies the floor, builds $\bs R$ with \code{rt\_of\_rv}
  (\code{src/minitensor.jl}) and returns the ideal element
  $\bs F_M^{-1} \bs Y$ together with $\bs F_M$ and $\bs F_M^{-1}$.
\item \code{evaluate\_element!}: with a metric field, the ideal element
  is built by the function above, the constitutive model is evaluated at
  $\bs F_U$, and the stress and the tangent are pulled back with
  \eqref{eq:pullback} by \code{pull\_back\_tangent}. Without a metric
  field the code path is unchanged.
\item Input validation knows the new keys (\code{metric field},
  \code{sizes}, \code{rotation vector}); the reference documentation is in
  \code{docs/src/reference/model.md}. The sampled sizes and rotation
  vector are written as the nodal variables \code{size\_1..3} and
  \code{rotation\_1..3}.
\end{itemize}

\section{Tests}
\label{sec:tests}

\code{test/smoothing-metric.jl} (test 126) checks:
\begin{enumerate}[itemsep=2pt]
\item the ideal element and the factor: $\bs F_M \bs F_M^{-1} = \bs I$,
  $\bs F_M^T \bs F_M = \bs M$ with and without rotation, unit edges of
  $\bs F_M \bs X$, the ideal volume, agreement with the isotropic size
  field for equal sizes, the floor (volume of the original element
  recovered, size ratios preserved), evaluation at the reference centroid
  and at the supplied time, and the input errors;
\item the orientation sensitivity on one element with sizes
  $(0.25, 1, 1)$ rotated by $0.4$ rad about $z$: the ideal element has
  zero energy, the same shape rotated by a quarter turn about $z$ has
  energy above one while the isotropic energy on $\bs F$ is zero for
  both, a rotation about the axis of transverse isotropy of the metric
  keeps the energy at zero, and $\bs Q \bs F_M$ gives the same energy as
  $\bs F_M$;
\item the assembled force and stiffness on the cube mesh
  (\code{examples/ems/cube}, 7677 elements) under a metric with all six
  expressions varying in space: central finite differences of the energy
  match the force to $10^{-5}$ relative and a finite-difference
  directional derivative of the force matches the stiffness to
  $10^{-5}$ relative;
\item a smoothing run on the cube with a constant metric at
  $45^\circ$ in the $x$-$y$ plane (sizes $0.5H$, $1.5H$, $1.5H$ with $H$
  the mean edge length): the energy decreases, the mean deviation of the
  edge lengths measured in the metric from unity decreases, and no
  element inverts.
\end{enumerate}
The existing smoothing tests (51, 119) and the input-validation sweep of
the examples (122) pass unchanged.

\section{Examples}
\label{sec:examples}

\paragraph{Cube, constant metric at $45^\circ$}
(\code{examples/ems/cube/cube-metric.yaml}). Sizes $0.05$, $0.1414$,
$0.1414$ (the volume of the mesh size $0.1$), first principal direction
at $45^\circ$ in the $x$-$y$ plane, unrestricted rule, nodes fixed in the
normal direction on the faces, L-BFGS. Over the edges of the mesh the mean
of $\lvert\, \lVert \bs F_M \bs e \rVert - 1 \rvert$ goes from $0.44$ to
$0.31$, the mean edge component along the small direction from $0.055$ to
$0.042$ (target $0.05$), and the mean component across it from $0.085$ to
$0.098$. The number of elements is fixed, so the metric cannot be met
everywhere; the elements rotate and stretch toward it.

\paragraph{Tube, graded axial size}
(\code{examples/ems/tube/tube-metric.yaml}). The rotation vector
$(0, -\pi/2, 0)$ carries the first principal direction onto the axis of
the tube; the axial size grows linearly from $0.06$ at the bottom to
$0.18$ at the top while the radial and circumferential sizes stay at the
mesh size $0.12$, restricted rule, nodes sliding on the four analytic
surfaces. The mean axial size equals the mesh size, so the grading can be
realized by sliding without a change of topology. The mean of
$\lvert\, \lVert \bs F_M \bs e \rVert - 1 \rvert$ over the edges goes
from $0.18$ to $0.12$, and the mean length of the axial edges (edges with
more than $70\%$ of their length along $z$) in the four quarters of the
height, from bottom to top, goes from $0.105$, $0.109$, $0.109$, $0.105$
to $0.075$, $0.110$, $0.127$, $0.144$, against targets of $0.075$,
$0.105$, $0.135$, $0.165$ at the quarter centers. The largest nodal motion
is $0.28$, about two element sizes, with no element inversion.

\paragraph{What smoothing cannot do.} A target that changes the number
of elements across the wall or around the circumference, such as a
radial size of a quarter of the wall thickness on a mesh with one element
through the wall, is not reachable by smoothing: the mesh barely moves and
the misfit does not decrease. Such targets require the topological
operations of the adaptivity note.

\section{Next steps}
\label{sec:next}

\begin{enumerate}[itemsep=2pt]
\item Nodal metrics carried by the mesh (six nodal variables: the
  logarithms of the sizes and the rotation vector), interpolated to the
  centroid with constant weights, and re-sampled from the field at phase
  boundaries.
\item Metrics given as full tensors: recovery of sizes and rotation vector
  by eigendecomposition with a consistent frame (eigenvector signs by
  continuity, a fixed in-plane phase for repeated sizes), with the
  log-Euclidean interpolation of the tensor as the fallback.
\item Coupling with the topological operations of the adaptivity note, so
  that targets the smoother cannot reach are met by refinement and
  coarsening in the metric.
\end{enumerate}

\end{document}
